\documentclass[11pt]{article}
\usepackage{reusv}
\usepackage{booktabs}
\usepackage{amsmath,amssymb}
\usepackage{siunitx}
\usepackage{graphicx}
\usepackage{multirow}
\usepackage{xcolor}
\usepackage{listings}

% ── Python 코드 스타일 ──
\definecolor{codebg}{HTML}{F5F5F0}
\definecolor{codeframe}{HTML}{CCCCBB}
\definecolor{codekw}{HTML}{0060A0}
\definecolor{codecomment}{HTML}{60802A}
\definecolor{codestring}{HTML}{B03030}
\lstdefinestyle{pythonstyle}{
  language=Python,
  backgroundcolor=\color{codebg},
  frame=single,
  rulecolor=\color{codeframe},
  framesep=6pt,
  xleftmargin=8pt,
  xrightmargin=8pt,
  basicstyle=\small\ttfamily,
  keywordstyle=\color{codekw}\bfseries,
  commentstyle=\color{codecomment}\itshape,
  stringstyle=\color{codestring},
  showstringspaces=false,
  breaklines=true,
  tabsize=4,
  numbers=left,
  numberstyle=\tiny\color{gray},
  numbersep=6pt,
  aboveskip=10pt,
  belowskip=6pt,
}

\setborderthickness{0.5pt}
\setbordermargin{1cm}
\setbordercolor{black}

\slimtitle{BLADE-SCP 수치 병목 개선 방법론}

\begin{document}

\drawborder
\thispagestyle{firstpage}

\lightertitle{BLADE-SCP 수치 병목 개선 방법론}

{\normalsize\textit{Research Note No.~017 $|$ bezier-orbit-transfer-scp}}\par\medskip

\def\lighterdate{March 12, 2026}
\lighterauthor{Suwon Lee$^{\dagger}$}

\lighteraddress{$^\dagger$}{Future Mobility Control Lab, Kookmin University, Seoul, Korea}
\lighteremail{suwon.lee@kookmin.ac.kr}

%% ── 초록 ─────────────────────────────────────────────────────
\begin{abstract}
BLADE 기반 SCP 궤도전이 최적화 솔버
\texttt{solve\_blade\_scp}의 수치 병목을 분석하고
4가지 최적화 기법을 적용하여 316\,ms에서 175\,ms로
1.81배 속도 향상을 달성하였다.
가장 큰 기여는 CVXPY의 DPP 캐싱 전략으로,
반복당 canonicalization 오버헤드를
57\,ms에서 3\,ms로 19배 절감하였다.
차수 올림의 closed-form 행렬화,
\texttt{bern\_product}의 einsum 퓨전,
Horner form 시도의 실패 분석을 함께 기술하여,
Bernstein 대수 프레임워크에서의
성능 최적화 방법론을 정리한다.
\end{abstract}

%% ══════════════════════════════════════════════════════════════
\section{도입}
%% ══════════════════════════════════════════════════════════════

보고서~011\cite{report011}에서 드리프트 적분의 벡터화를 통해
SCP 솔버의 실용적 속도를 확보하였다.
이후 BLADE 프레임워크\cite{efc005, report016}로 확장하면서,
\texttt{solve\_blade\_scp} 함수 전체에 대한 체계적 프로파일링을 수행하여
새로운 병목을 발견하고 이를 해결하였다.

본 보고서는 4가지 최적화 기법의 방법론과 성능 효과를 정리한다.
Bernstein 기저 기반 SCP 솔버의 일반적인 모듈 개발 시
참고할 수 있는 설계 지침을 제시하는 것이 목적이다.

최적화 전 프로파일링 결과, \texttt{solve\_blade\_scp}의 총 실행 시간
316\,ms (7회 SCP 반복)는 표~\ref{tab:before}과 같이 분해된다.

\begin{table}[h]
\centering
\caption{최적화 전 비용 분해 (316\,ms, 7 SCP iterations)}
\label{tab:before}
\begin{tabular}{lrr}
\toprule
구간 & 시간 [ms] & 비율 \\
\midrule
\texttt{algebraic\_drift} (gravity composition) & 147 & 46\% \\
CVXPY (canonicalization + solve) & 95 & 30\% \\
기타 ($\tau$ 갱신, BC 검사 등) & 58 & 18\% \\
\texttt{algebraic\_coupling} & 16 & 5\% \\
\midrule
\textbf{합계} & \textbf{316} & 100\% \\
\bottomrule
\end{tabular}
\end{table}

보고서~011\cite{report011}에서 벡터화한 \texttt{algebraic\_drift}가
여전히 최대 비중(46\%)을 차지하지만,
CVXPY의 canonicalization 오버헤드(30\%)가 두 번째로 큰 병목으로 부각되었다.
이하에서 각 최적화 기법을 효과 순으로 기술한다.


%% ══════════════════════════════════════════════════════════════
\section{CVXPY DPP 캐싱 전략}
\label{sec:dpp}
%% ══════════════════════════════════════════════════════════════

\subsection{문제: canonicalization 오버헤드}

CVXPY는 \texttt{Problem}을 풀기 전에 원뿔 표준형으로 변환하는
canonicalization 단계를 거친다.
이 과정에서 변수 구조, 제약 유형, 문제 크기를 분석하여 솔버에 전달할
행렬·벡터를 구성한다.

SCP 반복에서 매번 \texttt{cp.Problem}을 새로 구성하면,
\emph{문제 구조}는 동일한데 \emph{데이터}만 바뀌는 상황에서도
반복당 약 57\,ms의 canonicalization 비용이 발생한다.
나아가 \texttt{solve\_blade\_scp}를 여러 번 호출하는
outer loop(비행시간 탐색)에서는 첫 컴파일에 약 77\,ms가 추가로 소요된다.

\subsection{DPP(Disciplined Parameterized Programming) 원리}

CVXPY의 DPP 모드는 문제 구조를 1회만 컴파일하고,
이후에는 \texttt{cp.Parameter}의 값만 교체하여 재풀이할 수 있게 한다.
DPP 준수 조건은 다음과 같다:
\begin{enumerate}
  \item 변수(\texttt{cp.Variable})와 파라미터(\texttt{cp.Parameter})의
        곱·합이 DCP(Disciplined Convex Programming) 규칙을 따를 것
  \item 파라미터가 목적함수·제약의 \emph{계수}로만 나타날 것
        (파라미터 간 곱은 허용되지 않음)
\end{enumerate}

\subsection{적용 방법}

SCP 반복에서 변하는 데이터를 \texttt{cp.Parameter}로 선언한다:

\begin{lstlisting}[style=pythonstyle]
# Data that changes each iteration -> Parameter
A_eq = cp.Parameter((m_eq, n_var))  # equality LHS
b_eq = cp.Parameter(m_eq)           # equality RHS
J_ta = cp.Parameter((n_ta, n_var))  # thrust alloc Jacobian
lb   = cp.Parameter(n_var)          # trust region lower
ub   = cp.Parameter(n_var)          # trust region upper
w_bc = cp.Parameter(nonneg=True)    # BC penalty weight
\end{lstlisting}

문제 구조(변수, 목적함수, 제약 원뿔)는 1회만 구성하며,
\texttt{t\_f}와 \texttt{u\_max}도 \texttt{cp.Parameter}로 만들어
동일한 $(K, n, \texttt{ta\_free}, \texttt{trust\_region})$ 구조를 갖는
다른 문제에서도 캐싱된 문제를 재사용한다.

\begin{lstlisting}[style=pythonstyle]
# Module-level caching
_socp_cache: dict[tuple, _CachedSOCP] = {}

def get_or_create_socp(K, n, ta_free, trust_region):
    key = (K, n, ta_free, trust_region)
    if key not in _socp_cache:
        _socp_cache[key] = _build_socp(K, n, ta_free,
                                       trust_region)
    return _socp_cache[key]
\end{lstlisting}

\subsection{Trust region의 수치적 주의사항}

Box constraint로 trust region을 구현할 때,
비활성 변수에 대해 $\pm 10^{15}$ 같은 극단적 범위를 설정하면
내점법 솔버(Clarabel)의 complementary slackness 조건에서
수치 문제가 발생할 수 있다.
이를 방지하기 위해 $\pm (20 \times r_{\mathrm{trust}})$ 정도의
합리적 범위를 사용한다.

\subsection{효과}

\begin{itemize}
  \item 반복당 canonicalization: 57\,ms $\to$ 3\,ms (\textbf{19배} 절감)
  \item 호출 간 컴파일: 77\,ms $\to$ 0\,ms (캐시 적중 시)
  \item \texttt{warm\_start=True}로 이전 해를 시작점으로 활용하여
        솔버 반복도 감소
\end{itemize}


%% ══════════════════════════════════════════════════════════════
\section{차수 올림의 Closed-Form 행렬화}
\label{sec:degree_elevate}
%% ══════════════════════════════════════════════════════════════

\subsection{문제: 반복 루프에 의한 차수 올림}

$B_N \to B_R$ 차수 올림(\texttt{degree\_elevate})의 기존 구현은
$R - N$회의 1차 올림을 반복하며, 매 반복마다 새 배열을 할당한다.
BLADE-SCP에서 이 함수는 1,512회 호출되어 총 0.021\,s를 소요하였다.

\subsection{Closed-form 변환 행렬}

$N$차 Bernstein 제어점 $\mathbf{p} \in \mathbb{R}^{N+1}$을
$R$차 제어점 $\mathbf{p}' \in \mathbb{R}^{R+1}$로 직접 변환하는
행렬 $E_{N \to R}$을 도출한다:
\begin{equation}
  E_{N \to R}[k, j] = \frac{\binom{N}{j}\binom{R-N}{k-j}}{\binom{R}{k}},
  \quad 0 \le k \le R, \quad \max(0, k\!-\!R\!+\!N) \le j \le \min(N, k)
  \label{eq:elevate_matrix}
\end{equation}
이 행렬은 $(N, R)$ 쌍에만 의존하므로 \texttt{@functools.lru\_cache}로
캐싱하며, 차수 올림은 단일 행렬-벡터 곱으로 수행된다:
\begin{equation}
  \mathbf{p}' = E_{N \to R} \, \mathbf{p}
\end{equation}

\subsection{효과}

\begin{itemize}
  \item 함수 호출당: 14\,$\mu$s $\to$ 0.5\,$\mu$s (\textbf{28배})
  \item SCP 전체 기여: 약 10\,ms 절감
\end{itemize}


%% ══════════════════════════════════════════════════════════════
\section{\texttt{bern\_product} einsum 최적화}
\label{sec:einsum}
%% ══════════════════════════════════════════════════════════════

\subsection{문제: 불필요한 배열 할당과 임시 배열}

\texttt{bern\_product}의 벡터화 구현(보고서~011\cite{report011})은
다음 두 가지 비효율을 포함하고 있었다:

\begin{enumerate}
  \item \texttt{np.append(q, 0.0)}: 매 호출마다 새 배열 할당
        (7,812회 호출 시 누적 비용 유의미)
  \item 3-way broadcast
        \texttt{W * p[newaxis,:] * Q\_shifted}:
        중간 임시 배열 2개 생성
\end{enumerate}

\subsection{개선}

\paragraph{배열 할당 제거.}
\texttt{np.append} 대신 \texttt{np.minimum}에 의한 인덱스 클리핑과
유효 범위 마스크를 사용하여 원본 배열에서 직접 인덱싱한다.

\paragraph{einsum 퓨전.}
3-way broadcast를 \texttt{np.einsum}으로 퓨전하여
임시 배열 생성을 1개로 줄인다:
\begin{lstlisting}[style=pythonstyle]
# Before: 2 temporary arrays
r = np.sum(W * p[np.newaxis,:] * Q_shifted, axis=1)

# After: einsum fusion, 1 temporary array
Wp = W * p[np.newaxis, :]  # 1 temp array
r = np.einsum("kj,kj->k", Wp, Q_shifted)
\end{lstlisting}

\subsection{효과}

함수 호출 시간 약 10\% 절감.
\texttt{bern\_product}가 \texttt{bern\_compose} 내에서
수십 회 호출되므로, 절대적 기여는 수 ms 수준이다.


%% ══════════════════════════════════════════════════════════════
\section{Horner Form 시도와 교훈}
\label{sec:horner}
%% ══════════════════════════════════════════════════════════════

\subsection{시도: de Casteljau 방식의 Horner 변환}

\texttt{bern\_compose}는 합성 $f \circ g$의 Bernstein 제어점을
거듭제곱 전개로 계산한다.
$K$차 합성에서 순차 거듭제곱 방식은 $2K+1$회의
\texttt{bern\_product}를 호출하며,
이를 Horner form으로 변환하면 $K$회로 줄일 수 있을 것으로 기대하였다.

구체적으로, de Casteljau 알고리즘을 적용하여
Bernstein 계수를 재귀적으로 결합하는 방식을 시도하였다.

\subsection{결과: 오히려 느려짐}

\begin{table}[h]
\centering
\caption{Horner form 시도 결과 ($K = 12$)}
\label{tab:horner}
\begin{tabular}{lrrr}
\toprule
방법 & 시간 [ms] & 배속 & \texttt{bern\_product} 호출 수 \\
\midrule
순차 거듭제곱 (기존) & 0.185 & 1.00$\times$ & $3K+1 = 37$ \\
de Casteljau (Horner) & 0.482 & 0.38$\times$ & $K(K+1)/2 = 78^*$ \\
\bottomrule
\multicolumn{4}{l}{\footnotesize $^*$ 레벨당 $(K-r)$개 쌍, 총 $\sum_{r=0}^{K-1}(K-r) = K(K+1)/2$} \\
\end{tabular}
\end{table}

\subsection{원인 분석}

de Casteljau 알고리즘은 레벨 $r$에서 $(K-r)$개의 인접 쌍에 대해
각각 \texttt{bern\_product}를 호출한다.
$K=12$일 때 총 호출 수는
$\sum_{r=0}^{11}(12-r) = 78$회로,
기존 방식의 37회보다 2배 이상 많다.

초기 레벨에서 다루는 배열의 차수는 작아
개별 \texttt{bern\_product} 호출은 가볍지만,
Python 함수 호출 오버헤드(함수 디스패치, 인자 패킹, 캐시 조회 등)가
호출당 약 2--5\,$\mu$s씩 누적되어 총 비용이 증가한다.

\subsection{설계 지침}

Bernstein 대수 연산에서 성능 최적화 시
다음의 우선순위를 적용해야 한다:

\begin{enumerate}
  \item \textbf{함수 호출 횟수 최소화}:
    Python 레벨에서는 배열 크기보다 호출 횟수가 지배적 병목이다.
    알고리즘의 이론적 연산 복잡도보다
    실제 함수 호출 횟수를 기준으로 평가해야 한다.
  \item \textbf{호출당 연산량 극대화}:
    NumPy의 벡터화 이점을 최대한 활용하려면,
    작은 배열에 대한 다수의 호출보다
    큰 배열에 대한 소수의 호출이 효율적이다.
  \item \textbf{캐싱은 호출 횟수 감소에 준하는 효과}:
    \texttt{@lru\_cache}로 동일 입력의 재계산을 제거하면
    실질적 호출 횟수가 감소한다.
\end{enumerate}


%% ══════════════════════════════════════════════════════════════
\section{종합 성능 분석}
\label{sec:results}
%% ══════════════════════════════════════════════════════════════

\subsection{최적화 후 비용 분해}

표~\ref{tab:after}은 최적화 후의 비용 분해를 정리한다.
``핫(hot)'' 호출은 CVXPY 캐시가 적중한 경우,
``콜드(cold)'' 호출은 첫 컴파일이 포함된 경우이다.

\begin{table}[h]
\centering
\caption{최적화 후 비용 분해}
\label{tab:after}
\begin{tabular}{lrrr}
\toprule
구간 & 핫 [ms] & 콜드 [ms] & 최적화 전 [ms] \\
\midrule
CVXPY (canon.\ + solve) & 21 & 95 & 95 \\
\texttt{algebraic\_drift} & 112 & 112 & 147 \\
기타 & 42 & 76 & 74 \\
\midrule
\textbf{합계} & \textbf{175} & \textbf{283} & \textbf{316} \\
\bottomrule
\end{tabular}
\end{table}

\subsection{기법별 기여 요약}

표~\ref{tab:summary}에 각 기법의 개별 기여를 정리한다.

\begin{table}[h]
\centering
\caption{최적화 기법별 효과 요약}
\label{tab:summary}
\begin{tabular}{lccc}
\toprule
기법 & 개선 전 & 개선 후 & 배속 \\
\midrule
CVXPY DPP 캐싱 (반복당) & 57\,ms & 3\,ms & 19$\times$ \\
CVXPY DPP 캐싱 (호출 간) & 77\,ms & 0\,ms & --- \\
\texttt{degree\_elevate} closed-form & 14\,$\mu$s/call & 0.5\,$\mu$s/call & 28$\times$ \\
\texttt{bern\_product} einsum & --- & --- & $\sim$10\% \\
\texttt{bern\_compose} Horner & 0.185\,ms & 0.482\,ms & 0.38$\times$ \\
\midrule
\textbf{전체} (\texttt{solve\_blade\_scp}) & \textbf{316\,ms} & \textbf{175\,ms} & \textbf{1.81$\times$} \\
\bottomrule
\end{tabular}
\end{table}

\subsection{병목의 이동}

최적화 전에는 CVXPY와 드리프트 계산이
비슷한 비중(30\% vs 46\%)이었으나,
DPP 캐싱 후 CVXPY 비용이 크게 줄어
드리프트 계산이 전체의 64\%를 차지하게 되었다.
향후 추가 최적화는 Bernstein 곱/합성 연산의
C 확장 또는 JAX/Numba JIT 컴파일 방향이 유망하다.


%% ══════════════════════════════════════════════════════════════
\section{결론}
%% ══════════════════════════════════════════════════════════════

BLADE-SCP 솔버의 수치 병목을 체계적으로 분석하고
4가지 최적화 기법을 적용하여 1.81배 속도 향상을 달성하였다.
주요 결론은 다음과 같다.

\begin{enumerate}
  \item \textbf{CVXPY DPP 캐싱}이 가장 큰 효과를 보였다.
    \texttt{cp.Parameter}를 활용하여 문제 구조를 1회만 컴파일하고,
    모듈 레벨 캐싱으로 호출 간 재사용을 달성하면
    반복당 57\,ms의 canonicalization 오버헤드를 3\,ms로 줄일 수 있다.
    SCP처럼 동일 구조를 반복 풀이하는 알고리즘에서는
    DPP 설계가 필수적이다.

  \item \textbf{Closed-form 차수 올림 행렬}은
    이항계수 비로 직접 구성하여 28배의 가속을 얻었다.
    $(N, R)$ 쌍별 캐싱과 결합하면 반복 호출에서
    사실상 배열 조회 비용만 발생한다.

  \item \textbf{einsum 퓨전}은 임시 배열 생성을 줄여 약 10\% 개선을 주었다.
    NumPy 연산에서 중간 배열 할당은 소규모 배열일수록
    상대적 오버헤드가 크다.

  \item \textbf{Horner form 변환}은 이론적 연산량 감소에도 불구하고
    함수 호출 횟수 증가로 오히려 느려졌다.
    Bernstein 대수에서는 알고리즘의 이론적 복잡도보다
    Python 레벨 함수 호출 횟수가 성능을 지배한다는 교훈을 남겼다.
\end{enumerate}


%% ── 참고문헌 ─────────────────────────────────────────────────
\bibliographystyle{plain}
\bibliography{references}

\end{document}
